% RESEARCH_DRAFT_SHA256: e11a7baa4fad791db28fff6aab77cd13490d2786bd072671b790e9064abb5711
% Layout and source-reference wrappers expanded for the web renderer.
\documentclass[11pt]{article}
\usepackage[margin=25mm]{geometry}
\usepackage{fontspec,unicode-math}
\setmainfont{Latin Modern Roman}
\setmathfont{Latin Modern Math}
\usepackage{amsmath,amssymb,mathtools,xurl,hyperref}
\hypersetup{colorlinks=true,urlcolor=blue,linkcolor=blue}
\setlength{\parindent}{0pt}
\setlength{\parskip}{5pt}
\setlength{\emergencystretch}{3em}
\newcommand{\sref}[2]{\hyperlink{source:#1:#2}{[S#2]}}
\newcommand{\cataloguescope}[1]{\paragraph{Source question.}#1}
\newcommand{\floor}[1]{\left\lfloor#1\right\rfloor}
\begin{document}

\section*{Finite uniform CI representations: a rationalization route and its unproved step}
\noindent UnsolvedMath ID 20001197; AIM-COMPUTATION-0035. Research attempt dated 11 October 2026.
\paragraph{Result in scope.}This manuscript contains a source-specific attempt and complete proofs of the stated auxiliary calculations. It makes no new solution or novelty claim for the entire target.
\section{Original problem and preserved definition}
The following definition and source list are reproduced from the repository. Any earlier inspection dates and reported literature claims in this preserved text describe that source record, not additional verification by this attempt.
% BEGIN PRESERVED DEFINITION
\subsection{Definitions and mathematical statement}

\paragraph{Shared notation.}Standard number systems are denoted by $\mathbb N,\mathbb Z,\mathbb Q,\mathbb R,\mathbb C$; index ranges, computational models and specialized notation follow the cited problem source and the definitions below.


A conditional-independence (CI) structure records exactly which relations $X_A\perp X_B\mid X_C$ hold for disjoint index sets. It is probabilistically representable when it is the CI structure of some random vector. A finite uniform probability space gives each elementary outcome the same probability, while the random vector may identify multiple elementary outcomes.
\paragraph{Statement recorded in the supplied source.}
Conjecture 21. Every p-representable CI-structure can be p-represented by a random vector living on a finite probability space endowed with the uniform distribution. \href{https://aimath.org/WWN/compalgstat/compalgstat.pdf}{[S2]}
\paragraph{Scope qualification.}
The source uses p-representable for probabilistic representability; uniformity concerns elementary outcomes, not necessarily the distinct values of the random vector.
\paragraph{Source status.}
The supplied catalogue labels this record partially solved. Its source-status metadata was checked against the supplied snapshot on 2026-10-11. This import does not independently establish current open status or mathematical correctness.
\subsection{Short English statement}
Can every probabilistically realizable conditional-independence pattern be realized on a finite uniform probability space?
\subsection{Sources}
\begin{itemize}
\item[S1] Ulam AI and the UnsolvedMath Contributors, UnsolvedMath: A Curated Collection of Open Mathematics Problems, record 20001197 (AIM-COMPUTATION-0035). Catalogue curation and metadata: CC BY 4.0. The supplied source question is retained; definitions and the English summary are editorial additions.
\url{https://huggingface.co/datasets/ulamai/UnsolvedMath}.
\item[S2] AIM Workshop Problem Lists, original problem source linked by catalogue record AIM-COMPUTATION-0035
\url{https://aimath.org/WWN/compalgstat/compalgstat.pdf}.
\end{itemize}
% END PRESERVED DEFINITION
\subsection{Verbatim supplied catalogue question}
Conjecture 21. Every p-representable CI-structure can be p-represented by a random vector living on a finite probability space endowed with the uniform distribution.

\section{The exact representation requirement}
The source asks for exactly the same complete CI structure, not merely
for a distribution satisfying selected independences. Uniformity refers
to elementary outcomes, not to distinct vector values. We first make this
arithmetic requirement explicit, then try rational approximation within
a specified conditional-independence parametrization. The source does
not supply a finite-alphabet restriction on all initial representations;
our finite-alphabet argument does not establish such a reduction.

\paragraph{Finite uniform spaces and rational probability tables.}
A random vector on a uniform space of $N$ points has cell probabilities
$k_x/N$, where $k_x$ counts elementary outcomes with value $x$.
Conversely, if a finite table has rational nonnegative entries summing
to one, choose a common positive denominator $N$ and make $Np_x$
distinct elementary outcomes for each cell $x$, all of mass $1/N$.
Their total number is $N$ and the induced vector has exactly the
specified law. Consequently a rational finite table suffices for the
conjectured representation; equal positive masses at distinct vector
values are not necessary.

\paragraph{A rationalization theorem for a positive three-block chain.}
Let $X,Y,Z$ have fixed finite alphabets, a strictly positive joint table,
and $X\perp Z\mid Y$. Write
\[
p(x,y,z)=r_y a_{x\mid y}b_{z\mid y},
\]
where $r$ and each conditional vector $a,b$ are positive and sum to one.
Choose rational positive approximations to all but the last entry of
each vector and set the last entry to one minus their sum. Sufficiently
close approximations stay positive, and this polynomial
parametrization then produces rational tables converging to $p$, all
satisfying $X\perp Z\mid Y$ exactly.
To preserve the complete CI pattern, include every additional CI that
actually holds as a constraint on the parameter choice; the preceding
unrestricted approximation does not generally do that. Thus the theorem
proved by this argument is the following precisely limited one: if the
complete CI structure of $p$ has no additional nontrivial CIs beyond
those forced by $X\perp Z\mid Y$, then a rational table with that same
complete structure exists.
Here is the missing non-CI part of the proof. On fixed finite alphabets
there are finitely many disjoint-set CI statements. Each is expressed
by finitely many polynomial equations in table entries, for example
$p_{abc}p_c-p_{ac}p_{bc}=0$ for every assignment, with these symbols
meaning the corresponding marginal probabilities. At every absent CI,
at least one such polynomial is nonzero at $p$. Choose one witness for
each absent statement. By continuity, all the finitely many witnesses
remain nonzero in a sufficiently small neighborhood. The rational
parameter approximations can be taken in that neighborhood. All forced
CIs remain true by factorization and all absent CIs remain false by
the witnesses, proving the restricted claim. Trivial empty-set CIs
are automatically preserved.

\paragraph{General chart criterion and the exact unproved step.}
The same proof applies whenever a rational map, defined on an open
parameter neighborhood, takes that whole neighborhood into the asserted
CI locus and the probability simplex, takes one parameter point to the
given exact representation, has nonvanishing denominators there, and
has rational parameter points arbitrarily nearby. Rational outputs
preserve the asserted equations and finitely many nonzero witnesses
preserve every absent CI. This is a sufficient criterion, proved by
those two continuity arguments. It is not a proof that every CI locus
has such a chart.

For example, even the displayed chain approximation fails to preserve
an additional marginal independence unless its parameters satisfy the
additional polynomial exactly. General rational approximation in the
ambient simplex is still less adequate. The strictly positive normalized
set
\[
u^2=2v^2,\qquad u+v=1,\qquad u,v>0
\]
has a real point $v=1/(1+\sqrt2)$ but no rational point: rational
$u,v$ with $v>0$ would imply a rational square root of two. Its
irrationality follows by reducing $a^2=2b^2$ to coprime integers,
forcing both $a$ and $b$ even. This set is not asserted to be a CI
variety or a counterexample to the conjecture. It disproves only the
proposed general inference from rational defining polynomials and a real
point to rational points on the same locus.

\section{Audit and scope}
The AIM Conjecture 21 was inspected. The audit explicitly corrected an
initial overly broad chain-approximation idea: retaining one factorization
does not retain every additional CI of a special parameter point.
The final restricted theorem includes this essential hypothesis. It
also checks complete patterns, support positivity, common-denominator
refinement, and the missing finite-alphabet reduction. No numerical or
formal algebra computation was run. The general existence of a
rational chart or another exact rational representation is not proved;
no original conjectural progress or novelty is claimed.

\section{Process and contributions}
Mathematics and model-assisted adversarial self-audit: OpenAI gpt-6-astra, ultra reasoning, as dispatched for this research batch; exact serving revision not exposed. Codex (GPT-6; exact serving revision not exposed) materially assisted with orchestration and format review. No human mathematical authorship, mathematical edits, or independent human proof checking is asserted.
Contribution prompt A was applied to the full definition and supplied source row, then prompt B was applied adversarially to the derived results. The item-specific audit above records the scope checks and rejected extrapolations. Local assembly and structural checks do not establish mathematical correctness. The source catalogue status is preserved in TOP\_PROBLEM and is separate from this attempt status.
\paragraph{Integration review.} Codex (GPT-6; exact serving revision not exposed) reviewed the written arguments and expanded layout and source-reference wrappers. Repository metadata and the production web renderer were checked using Python 3.13 and JavaScript with MathJax 3.2.2. This model-assisted review is not independent human proof verification; no PDF compilation or formal proof checking is asserted.

\section{Completion Estimate}
\noindent\textbf{COMPLETION: 0\%}
The proved material is a rederivation of known or auxiliary facts rather than a substantiated new advance on the original target. Finite-uniform equivalence and a restricted rational-chart construction are proved, but arbitrary asserted-CI intersections and any finite-alphabet reduction remain unproved.
\section{References}
The complete original Sources subsection above is retained, including catalogue attribution and licence notices. The additional list records primary sources actually consulted for this attempt on 11 October 2026; uncited source links reproduced above were not all rechecked.
\begin{itemize}
\item American Institute of Mathematics, Computational Algebraic Statistics, prepared by Seth Sullivant, version 27 March 2004; the original numbered source question and its adjacent scope were inspected on 11 October 2026. \url{https://aimath.org/WWN/compalgstat/compalgstat.pdf}.
\end{itemize}
\section{Final status}
\textbf{Unresolved in this attempt.} Finite-uniform equivalence and a restricted rational-chart construction are proved, but arbitrary asserted-CI intersections and any finite-alphabet reduction remain unproved.
\end{document}
